% Compile from the course root with pdflatex; the editable master is this .tex file.
% Production: print double-sided, long-edge binding. Copy count equals enrolled
% students excluding cancelled exam sittings; do not request extra copies.
% Format: Garamond 11 pt, single spacing, 2 cm margins, black; section headings bold, no shading.
\documentclass[a4paper,11pt,twoside]{article}
\usepackage[T1]{fontenc}
\usepackage{ebgaramond}
\usepackage{amsmath,amssymb}
\usepackage[ebgaramond]{newtxmath}
\usepackage[margin=2cm,includefoot,footskip=8mm]{geometry}
\usepackage{enumitem,tabularx,fancyhdr,needspace}
\linespread{1}
\setlength{\parindent}{0pt}
\setlength{\parskip}{3pt}
\setlength{\abovedisplayskip}{5pt}
\setlength{\belowdisplayskip}{5pt}
\setlength{\abovedisplayshortskip}{3pt}
\setlength{\belowdisplayshortskip}{3pt}
\setlength{\emergencystretch}{1em}
\raggedbottom
\pagestyle{fancy}
\fancyhf{}
\renewcommand{\headrulewidth}{0pt}
\renewcommand{\footrulewidth}{0pt}
\fancyfoot[C]{\normalfont\thepage}
\pdfcatalog{/ViewerPreferences << /Duplex /DuplexFlipLongEdge >>}
\newcommand{\E}{\operatorname{E}}
\newcommand{\Var}{\operatorname{Var}}
\newcommand{\Cov}{\operatorname{Cov}}
\newcommand{\Corr}{\operatorname{Corr}}
\newcommand{\mcoptions}[4]{\par\nobreak\smallskip
  \begin{tabularx}{\linewidth}{@{}X X@{}}
  (a) #1 & (b) #2 \\
  (c) #3 & (d) #4
  \end{tabularx}}
\makeatletter
\renewcommand{\section}{\@startsection{section}{1}{\z@}%
  {6pt}{3pt}{\normalfont\bfseries\fontsize{12}{14}\selectfont}}
\makeatother
\newsavebox{\ExamOpeningBox}
\newlength{\ExamOpeningHeight}
\newcommand{\ExamOpening}{%
  \begin{minipage}[b]{\textwidth}
  \normalfont\normalsize
  \linespread{1}\selectfont
  \setlength{\parindent}{0pt}
  \setlength{\parskip}{3pt}
  \begin{center}
  {\fontsize{12}{14}\selectfont Econometrics (25117): Midterm 1, 2026--2027}\\
  Duration: 180 minutes\\
  \normalsize Monday, 9 November 2026
  \end{center}
  Name and surname: \rule{8.4cm}{0.4pt}\hfill UPF ID: U\makebox[4mm]{\_}\makebox[4mm]{\_}\makebox[4mm]{\_}\makebox[4mm]{\_}\makebox[4mm]{\_}\makebox[4mm]{\_}\par
  Answer all questions. Questions 1--8 have one correct option each: 2.5 points for a
  correct answer, 0 for a wrong or blank answer. Questions 9--15 are worth 80 points;
  show calculations and briefly explain your reasoning on separate answer paper.
  The total is 100 points; the grade out of 10 is the total divided by 10.
  For inference, use the large-sample standard normal approximation. The two-sided
  critical values are 1.96 at 5\% and 2.576 at 1\%.
  Exact fractions, radicals or sensible rounding are acceptable. All numerical examples are
  hypothetical.\par
  \vspace{2pt}\noindent\leaders\hbox{.\kern2pt}\hfill\kern0pt\par\vspace{2pt}
  \end{minipage}}

\begin{document}
\sbox{\ExamOpeningBox}{\ExamOpening}
\setlength{\ExamOpeningHeight}{\dimexpr\ht\ExamOpeningBox+\dp\ExamOpeningBox+5pt\relax}
\newgeometry{left=2cm,right=2cm,top=2cm,bottom=2cm,includehead,includefoot,headheight=\ExamOpeningHeight,headsep=3mm,footskip=8mm}
\setlength{\headwidth}{\textwidth}
\fancyhead[C]{\usebox{\ExamOpeningBox}}

\section*{Multiple-choice questions}
\begin{enumerate}[label=\arabic*.,leftmargin=*,itemsep=5pt,topsep=4pt,parsep=0pt]
\item Let $Y_1,\ldots,Y_n$ be i.i.d. with mean $\mu$ and variance
$0<\sigma^2<\infty$. Consider $\widetilde\mu_n=Y_1$, which uses only the first
observation, regardless of $n$. This estimator is:
\mcoptions{unbiased and consistent.}{unbiased and inconsistent.}
{biased and consistent.}{biased and inconsistent.}

\item Which statement always holds when the required moments exist?
\mcoptions{Zero covariance implies independence.}
{Zero covariance implies $\E[Y\mid X]=\E[Y]$.}
{Independence implies zero covariance.}
{Perfect positive correlation implies independence.}

\item For i.i.d. observations with mean $\mu$ and variance
$0<\sigma^2<\infty$, the central limit theorem states that, as $n\to\infty$:
\mcoptions{$\bar Y$ has exactly a normal distribution for every $n$.}
{$\sqrt{n}(\bar Y-\mu)/\sigma\xrightarrow{d}N(0,1)$.}
{$(\bar Y-\mu)/\sigma\xrightarrow{d}N(0,1)$.}
{$\bar Y\xrightarrow{p}0$, whatever the value of $\mu$.}

\item A Type I error occurs when a test:
\mcoptions{rejects a true null hypothesis.}
{does not reject a false null hypothesis.}
{rejects a false null hypothesis.}
{does not reject a true null hypothesis.}

\item A large-sample two-sided test uses an estimate $\bar y=3$ and
$SE(\bar Y)=1$. The null $H_0:\mu_Y=0$ is rejected at 5\%. For any
hypothesized value $\mu_0$, reject when $|t|>1.96$. What is the smallest
$\mu_0$ for which $H_0:\mu_Y=\mu_0$ would not be rejected?
\mcoptions{$0$.}{$1.96$.}{$4.96$.}{$1.04$.}

\item In an OLS regression with an intercept and a nonconstant regressor,
which statement is necessarily true, apart from numerical rounding?
\mcoptions{$\sum_i\widehat U_i=0$ proves that $\E[U_i\mid X_i]=0$.}
{$\sum_i\widehat U_i=0$ and $\sum_i X_i\widehat U_i=0$, even if the regressor is endogenous.}
{$U_i=\widehat U_i$ for every observation.}
{Every individual residual equals zero.}

\item For i.i.d. observations with a finite, positive population variance,
the standard deviation of the sample mean is $\sigma_Y/\sqrt n$.
If the sample size increases from 100 to 400 while the population distribution
stays the same, the standard deviation of the sample mean:
\mcoptions{is divided by four.}{is multiplied by two.}
{is unchanged.}{is divided by two.}

\item For simple OLS with an intercept, $n>2$ and $SSR>0$, which statement
correctly distinguishes the standard error of the regression ($SER$) from
the root mean squared residual ($RMSE$)?
\mcoptions{$SER=\sqrt{SSR/n}$; $RMSE=\sqrt{SSR/(n-2)}$.}
{$SER=RMSE$ for every finite sample.}
{$SER=\sqrt{SSR/(n-2)}$; $RMSE=\sqrt{SSR/n}$.}
{Both equal $\sqrt{SSR/(n-1)}$.}

\end{enumerate}
\Needspace{7.5cm}
\section*{Long-answer questions}
\begin{enumerate}[label=\arabic*.,leftmargin=*,itemsep=5pt,topsep=4pt,parsep=0pt,start=9]
\Needspace{6.5cm}
\item Joint probabilities and dependence (10 points).
Let $D=1$ indicate attendance at an optional career workshop and $Y=1$ indicate
employment after graduation. Their population joint distribution is:
\begin{center}
\begin{tabular}{c|cc}
 & $D=0$ & $D=1$ \\ \hline
$Y=0$ & 0.36 & 0.24 \\
$Y=1$ & 0.04 & 0.36
\end{tabular}
\end{center}
\begin{enumerate}[label=(\alph*),leftmargin=*,itemsep=2pt,topsep=2pt]
\item Calculate $\Pr(D=1)$ and $\Pr(Y=1)$. (2 points)
\item Calculate $\E[Y\mid D=0]$ and $\E[Y\mid D=1]$, then use the law of
iterated expectations to recover $\E[Y]$. (3 points)
\item Calculate $\Cov(D,Y)$ and $\Corr(D,Y)$. (3 points)
\item Are $D$ and $Y$ independent? Justify your answer with a probability
calculation. (2 points)
\end{enumerate}

\item Comparing two means (10 points).
An experiment randomly assigns students to a tutoring group ($T$) or a control
group ($C$). Treat observations as i.i.d. within each group and independent across groups.
The outcome is a test score. The sample summaries are:
\begin{center}
\begin{tabular}{lccc}
Group & Sample size & Sample mean & Sample standard deviation \\ \hline
Tutoring & 100 & 74 & 12 \\
Control & 100 & 68 & 16
\end{tabular}
\end{center}
Let $\Delta=\mu_T-\mu_C$ be the population difference in mean scores.
\begin{enumerate}[label=(\alph*),leftmargin=*,itemsep=2pt,topsep=2pt]
\item Estimate $\Delta$ and calculate its estimated standard error. (2 points)
\item The claim is that tutoring raises the population mean by exactly two
points. State a null and a two-sided alternative, and calculate the test
statistic. (2 points)
\item Use the supplied critical values to decide whether to reject the null
at the 5\% and 1\% significance levels. Show both comparisons. (2 points)
\item Construct a 95\% confidence interval for $\Delta$. Does it contain the
null value from (b)? Explain how this relates to your test decision. (2 points)
\item Plan a new study with $m$ students in each group. Using 12 and 16 as the
planning population standard deviations, calculate the smallest integer $m$
for a planned 95\% confidence-interval half-width of at most 2 score points.
Also report the total planned sample size. (2 points)
\end{enumerate}

\Needspace{6.5cm}
\item A small library population (12 points).
The following table describes a population of eight students. $G$ records each
student's programme, A or B, and $Y$ is the number of books borrowed last week.
The entries are numbers of students, not probabilities. Select one of these
eight students uniformly at random.
\begin{center}
\begin{tabular}{c|cc}
Books borrowed, $Y$ & Programme A & Programme B \\ \hline
0 & 1 & 0 \\
2 & 2 & 4 \\
4 & 1 & 0
\end{tabular}
\end{center}
\begin{enumerate}[label=(\alph*),leftmargin=*,itemsep=2pt,topsep=2pt]
\item Give the marginal probability distribution of $Y$ and write its complete
cumulative distribution function $F_Y(y)=\Pr(Y\leq y)$. (3 points)
\item Calculate $\E[Y]$, $\E[Y^2]$ and the population variance $\Var(Y)$.
(3 points)
\item Calculate $\E[Y\mid G=\mathrm A]$ and $\E[Y\mid G=\mathrm B]$. Use
these conditional means and the programme probabilities to recover $\E[Y]$
by the law of iterated expectations. (3 points)
\item Are $G$ and $Y$ independent? Justify your answer with a conditional
probability calculation. Explain what this example shows about using equal
conditional means to infer independence. (3 points)
\end{enumerate}

\Needspace{6.5cm}
\item From sample summaries to inference (12 points).
An i.i.d. sample of 100 students has assessment scores $y_i$, measured in
points. The following raw sums are available:
\[
n=100,\qquad \sum_{i=1}^{100}y_i=5200,\qquad
\sum_{i=1}^{100}y_i^2=280300.
\]
Assume that the large-sample standard normal approximation is adequate.
Use the two-sided critical values 1.96 at 5\% and 2.576 at 1\%; no
probability-table lookup is needed. You may use the algebraic identity
\[
\sum_i(y_i-\bar y)^2=\sum_i y_i^2-n\bar y^2.
\]
\begin{enumerate}[label=(\alph*),leftmargin=*,itemsep=2pt,topsep=2pt]
\item Calculate the sample mean, the centered sum of squares, the unbiased
sample variance, the sample standard deviation, and the estimated standard
error of the sample mean. State the units of the standard deviation and
the standard error. (3 points)
\item Test the claim that the population mean score equals 50 points against
a two-sided alternative. State both hypotheses, calculate the test statistic,
and give the decisions at 5\% and 1\%, showing both critical-value
comparisons. (3 points)
\item Calculate the 95\% confidence interval for the population mean.
Keep at least two decimal places in its endpoints. Does it contain 50?
Relate this answer to your 5\% test decision. (3 points)
\item Plan a new i.i.d. survey using 10 points as the planning population
standard deviation. Find the minimum integer sample size for a planned
95\% interval half-width of at most 0.98 points. Compare this sample size
with 100 and explain the change in precision. (3 points)
\end{enumerate}

\item Simple OLS: estimation, fit and interpretation (12 points).
Using $n=102$ independent observations, a researcher regresses test score $Y_i$
on weekly independent-study hours $X_i$, with an intercept:
$Y_i=\beta_0+\beta_1X_i+U_i$. The observed sample has
\[
\bar x=4,\quad \bar y=70,\quad
\sum_i(x_i-\bar x)^2=100,\quad
\sum_i(x_i-\bar x)(y_i-\bar y)=50,\quad
TSS_Y=\sum_i(y_i-\bar y)^2=425.
\]
\begin{enumerate}[label=(\alph*),leftmargin=*,itemsep=2pt,topsep=2pt]
\item Before substituting the sample summaries, write the least-squares objective for the model with an intercept,
differentiate with respect to both coefficients, and derive formulas for
$\widehat\beta_0$ and $\widehat\beta_1$. State the condition required for a unique
solution. (2 points)
\item Calculate $\widehat\beta_1$ and $\widehat\beta_0$. Write the fitted line
and interpret the slope in the units of this example. (2 points)
\item For a student with $x_i=6$ and $y_i=72$, calculate the fitted score and
the residual. (2 points)
\item Calculate the explained sum of squares ($ESS$), the residual sum of
squares ($SSR$), $R^2$ and $SER$. (2 points)
\item Calculate the sample correlation between study hours and test scores
using the supplied centered sums. Verify numerically that the square of this correlation equals
$R^2$. (2 points)
\item Study time is now recorded in minutes, $M_i=60X_i$. Calculate the
intercept and slope in the fitted regression of test score on minutes.
Report the numerical values of $R^2$ and $SER$ after this conversion.
(2 points)
\end{enumerate}


\Needspace{7cm}
\item One regressor without an intercept (12 points).
Use the same 102 observations and sample summaries given in Question 13.
Consider the model $Y_i=\beta_1X_i+U_i$, with no intercept. Derive the
estimator before substituting numbers. Exact fractions are acceptable
throughout; decimal division is not required.
\begin{enumerate}[label=(\alph*),leftmargin=*,itemsep=2pt,topsep=2pt]
\item Write the least-squares objective, differentiate it to obtain the
first-order condition, and derive a closed-form formula for
$\widehat\beta_1$. State the denominator condition required for a unique
solution. (4 points)
\item Calculate $\widehat\beta_1$ and
interpret the fitted slope in the units of this example. (4 points)
\item For the same student with $x_i=6$ and $y_i=72$ in Question 13,
calculate the fitted score and residual under this no-intercept model.
(2 points)
\item Use the first-order condition to evaluate
$\sum_i x_i\widehat u_i$. Calculate the ordinary residual sum
$\sum_i\widehat u_i$. Must this ordinary sum equal zero without an
intercept? Explain. (2 points)
\end{enumerate}

\Needspace{6.5cm}
\item Simple OLS from five observations (12 points).
The following table contains five hypothetical students. $X$ is weekly study
time in hours and $Y$ is a quiz score in points.
\begin{center}
\begin{tabular}{c|rrrrr}
Student $i$ & 1 & 2 & 3 & 4 & 5 \\ \hline
$x_i$ & 1 & 2 & 3 & 4 & 5 \\
$y_i$ & 4 & 4 & 7 & 8 & 12
\end{tabular}
\end{center}
Fit simple OLS with an intercept, $\widehat y_i=\widehat\beta_0+
\widehat\beta_1x_i$. All calculations can be completed by hand; give exact
fractions or radicals rather than decimal approximations when appropriate.
\begin{enumerate}[label=(\alph*),leftmargin=*,itemsep=2pt,topsep=2pt]
\item Calculate $\bar x$, $\bar y$, $S_{XX}=\sum_i(x_i-\bar x)^2$ and
$S_{XY}=\sum_i(x_i-\bar x)(y_i-\bar y)$. Use them to estimate the slope and
intercept. Write the fitted line and interpret the slope. (4 points)
\item Calculate the fitted score and residual for every student. Verify
numerically both $\sum_i\widehat u_i=0$ and
$\sum_i x_i\widehat u_i=0$. (3 points)
\item Calculate $TSS=\sum_i(y_i-\bar y)^2$,
$ESS=\sum_i(\widehat y_i-\bar y)^2$, $SSR=\sum_i\widehat u_i^2$,
$R^2$ and $SER$. Show that $TSS=ESS+SSR$. (3 points)
\item Study time is instead recorded in minutes, $M_i=60X_i$. Calculate the
intercept and slope of the same fitted line expressed in minutes. Report
$R^2$ and $SER$ in these units, briefly explaining whether they change.
(2 points)
\end{enumerate}
\end{enumerate}
\end{document}
